% ============================================================================
\section{Dispatch 层：\texttt{syscall.c} 实现}
\label{sec:syscall-dispatch}
% ============================================================================

\texttt{syscall.c} 是系统调用子系统的核心——持有 syscall table，
实现后端注册、handler 注册、和调用分发。

% ----------------------------------------------------------------------------
\subsection{syscall table}
% ----------------------------------------------------------------------------

\begin{figure}[H]
  \centering
  \begin{tikzpicture}[
    cell/.style={draw, minimum width=2.5cm, minimum height=0.65cm,
                 font=\ttfamily\small, anchor=north},
    idx/.style={font=\ttfamily\scriptsize\color{gray}, anchor=east},
    used/.style={cell, fill=green!10},
    empty/.style={cell, fill=gray!5},
    arrow/.style={-{Stealth[length=3pt]}, thick, blue!60!black},
  ]

  % 数组框
  \node[used]  (s0) at (0, 0)     {NULL};
  \node[used]  (s1) at (0, -0.65) {sys\_exit};
  \node[empty] (s2) at (0, -1.30) {NULL};
  \node[empty] (s3) at (0, -1.95) {NULL};
  \node[used]  (s4) at (0, -2.60) {sys\_write};
  \node[empty] (s5) at (0, -3.25) {\dots};
  \node[used]  (s45) at (0, -3.90) {sys\_brk};
  \node[empty] (s46) at (0, -4.55) {\dots};
  \node[empty] (s255) at (0, -5.20) {NULL};

  % 索引标签
  \node[idx] at (-1.5, 0)     {[0]};
  \node[idx] at (-1.5, -0.65) {[1]};
  \node[idx] at (-1.5, -1.30) {[2]};
  \node[idx] at (-1.5, -1.95) {[3]};
  \node[idx] at (-1.5, -2.60) {[4]};
  \node[idx] at (-1.5, -3.25) {};
  \node[idx] at (-1.5, -3.90) {[45]};
  \node[idx] at (-1.5, -4.55) {};
  \node[idx] at (-1.5, -5.20) {[255]};

  % EAX 指向
  \node[font=\small, anchor=east] (eax) at (-3.5, -2.60) {\textbf{EAX=4}};
  \draw[arrow] (eax.east) -- (s4.west);

  % 标题
  \node[font=\small\bfseries] at (0, 0.8) {\texttt{syscall\_table[SYSCALL\_MAX]}};

  % 说明
  \node[font=\scriptsize\itshape, text=gray, anchor=west, text width=3cm]
    at (1.8, -1.0) {预留 0 号不使用\\（检测 EAX 未初始化）};
  \node[font=\scriptsize\itshape, text=gray, anchor=west, text width=3cm]
    at (1.8, -4.0) {编号参考 Linux i386\\但只实现需要的子集};

\end{tikzpicture}

  \caption{syscall table：系统调用号直接索引函数指针数组}
  \label{fig:syscall-table}
\end{figure}

\codefile{src/kernel/core/syscall.c}
\begin{lstlisting}[style=cstyle]
static syscall_handler_t syscall_table[SYSCALL_MAX];

/* 当前注册的后端 */
static const syscall_ops_t* g_syscall_ops = NULL;
static int g_syscall_ready = 0;
\end{lstlisting}

\begin{whybox}
为什么用静态数组而非链表或哈希表？
\begin{enumerate}
  \item 系统调用号是小整数（通常 $< 256$），数组索引 \textbf{$O(1)$}——没有比这更快的了
  \item 无动态内存分配——启动早期 kmalloc 可能还没初始化
  \item 256 个指针只占 1~KiB，内存开销可忽略
  \item 未注册的 slot 为 NULL，dispatch 时一次判断即可
\end{enumerate}
\end{whybox}

% ----------------------------------------------------------------------------
\subsection{后端注册}
% ----------------------------------------------------------------------------

\begin{lstlisting}[style=cstyle]
void syscall_register_backend(const syscall_ops_t* ops) {
    g_syscall_ops = ops;
}

const char* syscall_backend_name(void) {
    return g_syscall_ops ? g_syscall_ops->name : NULL;
}

int syscall_is_ready(void) {
    return g_syscall_ready;
}
\end{lstlisting}

这与 \texttt{vma\_register\_backend()} 完全同一套模式——一行赋值即完成注册。

% ----------------------------------------------------------------------------
\subsection{handler 注册}
% ----------------------------------------------------------------------------

\begin{lstlisting}[style=cstyle]
int syscall_register(uint32_t nr, syscall_handler_t handler) {
    if (nr == 0 || nr >= SYSCALL_MAX) {
        printk("[syscall] register: invalid nr=%u\n", nr);
        return -1;
    }
    syscall_table[nr] = handler;
    return 0;
}
\end{lstlisting}

\begin{warnbox}
\texttt{nr == 0} 被拒绝——0 号保留为"未初始化 EAX"的安全网。
\texttt{nr >= SYSCALL\_MAX} 防止数组越界。
这两个检查在系统边界上是必要的，即使当前只有内核自己调用 \texttt{syscall\_register}。
\end{warnbox}

% ----------------------------------------------------------------------------
\subsection{syscall\_dispatch：调用分发}
\label{sec:syscall-dispatch-fn}
% ----------------------------------------------------------------------------

后端的汇编 stub 从寄存器提取参数，打包成 \texttt{syscall\_regs\_t}，
然后调用此函数。这是系统调用路径上最关键的 C 函数。

\begin{lstlisting}[style=cstyle]
#define ENOSYS 38

int32_t syscall_dispatch(const syscall_regs_t* regs) {
    if (!regs) {
        return -ENOSYS;
    }

    uint32_t nr = regs->nr;

    if (nr >= SYSCALL_MAX || !syscall_table[nr]) {
        printk("[syscall] unknown syscall nr=%u\n", nr);
        return -ENOSYS;
    }

    return syscall_table[nr](regs);
}
\end{lstlisting}

逻辑极简：
\begin{enumerate}
  \item 边界检查（\texttt{nr} 不超过 table 大小）
  \item NULL 检查（该 slot 是否注册了 handler）
  \item 调用 handler，返回结果
\end{enumerate}

\begin{thinkbox}
为什么返回 $-$ENOSYS（$-38$）而不是 $-1$？

\emph{提示：这是 POSIX 约定。\texttt{ENOSYS} 表示"Function not implemented"，
是 Linux 内核对未实现系统调用的标准返回值。
遵循这个约定，将来移植 GCC/libc 时用户态 \texttt{errno} 能正确设置。}
\end{thinkbox}
